\documentclass[10pt,a4paper]{amsart}
\usepackage[margin=19mm]{geometry}
\usepackage{amsmath,amssymb,mathtools}
\usepackage[scr=rsfs,cal=euler]{mathalfa}
\usepackage{xfrac}
\usepackage[unicode,hidelinks,bookmarks=false]{hyperref}
\newcommand{\ie}{i.e.\ }
\newcommand{\cf}{cf.\ }
\newcommand{\resp}{respectively\ }
\newcommand{\Subsection}{\subsection}
\providecommand{\nobreakdash}{\nobreak\hbox{-}}
\setlength{\emergencystretch}{2em}
\multlinegap0pt
\usepackage{tikz}
\usetikzlibrary{cd}
\usepackage{mathtools}
\usepackage{amssymb}
\usepackage{enumerate}
\usepackage[all]{xy}
\usepackage{xcolor}
\newtheorem{theorem}{Theorem}
\newtheorem{lemma}{Lemma}
\newcommand{\Mod}{\mathrm{Mod}}
\newcommand{\Qcoh}{\mathrm{QCoh}}
\newcommand{\ZZ}{\mathbf{Z}}
\newcommand{\cofib}{\mathrm{cofib}}
\newcommand{\id}{\mathrm{id}}
\newcommand{\jcomp}{{J\text{-}\mathrm{comp}}}
\newcommand{\qq}{/\!\!/}
\newcommand{\spec}{{\mathrm{Spec}}}
\newcommand{\spf}{\operatorname{Spf}}
\title{Derived algebras on formal stacks and prismatic gauges}
\author{Shubhankar Sahai}
\date{Source: arXiv:2602.08942v2 (CC BY 4.0)}
\begin{document}
\maketitle

\noindent\emph{Source and licence.} Derived algebras on formal stacks and prismatic gauges. Version arXiv:2602.08942v2 (CC BY 4.0), distributed by arXiv under the Creative Commons Attribution 4.0 International licence, http://creativecommons.org/licenses/by/4.0/, as recorded in the arXiv licence metadata for this exact version and shown on the abstract page https://arxiv.org/abs/2602.08942v2. Full licence notice: \texttt{licenses/arxiv-2602.08942-CC-BY-4.0.txt}.\par\smallskip

\noindent\textit{Editorial scope.} Counted unit: the article's lemma koszul (source0.tex lines 918--920). The article's complete original proof of it is delivered with it (source0.tex lines 921--966, one \texttt{proof} environment of 46 lines). No mathematics is written, repaired or re-derived: the statement and the proof are the article's own lines, in the article's own notation, and no retained line is substituted or re-flowed.  Retained uncounted as the source-local context the proof needs: lines 898--900.  Nothing was omitted from any retained span, so the package carries no omission record.  No retained line contains a citation, so the wrapper carries no bibliography and no bibliography entry.  The retained lines contain the article's own diagram code, which the wrapper typesets with the standard tikz packages; no image file is delivered and the package declares no asset dependency.  Editorial formatting only, in addition: an AMS \texttt{amsart} preamble that restates the article's own declarations and macros used by the retained lines, namely the environments \texttt{theorem}, \texttt{lemma} on separate counters, together with the standard packages those lines need.  The retained environments print in this document as Theorem 1, Lemma 1 in order of appearance, whereas the article prints the same items with the numbering of its own full document.  The complete native source of the article as distributed in the arXiv e-print for this exact version is bundled unchanged as \texttt{sources/194283/source0.tex}.
\begin{theorem}\label{thm: mod on spf}
    The pullback $(\eta_A^*)_{\Mod}\colon \Qcoh(\spec(A))\simeq \Mod_A\to  \Qcoh(\spf(A))$ restricts to an equivalence on $\Mod_A^\jcomp$
\end{theorem}

\begin{lemma}\label{lem: pro koszul}
    In the situation of Theorem \ref{thm: mod on spf}, consider the pro-system given by $\{A_{n}\otimes_A A_1\}_{n>0}$ and the constant pro-system $\{A_1\}_{n>0}$. Then $\{A_{n}\otimes_A A_1\}_{n>0}$ is pro-isomorphic to $\{A_1\}_{n>0}$.
\end{lemma}

\begin{proof}
Note that for any $n>0$ we have $A_n=A\qq f_1^n,\ldots, f^n_r$ is obtained by freely setting the elements $f^n_i=0$ for $ 1\leq i\leq r$. 


In particular if $M$ is any $A$-module, then the action of $x_i^n$ on $M\otimes_A A_n$ factors through $0.$
Since tensors are computed at the level of modules, we may also write the underlying as $\bigotimes_{i=1}^{i=r}\cofib(A\xrightarrow{x_i^n} A).$

Thus we compute 
\begin{align}\label{eq: formula for antimesa1}
A_n\otimes_A A_1
  &= (\bigotimes_{i=1}^{r}\cofib(A\xrightarrow{x_i^n} A))\otimes_A (\bigotimes_{i=1}^{r} A\qq x_i)\notag\\
  &= \bigotimes_{i=1}^{r}\cofib(A\qq x_i\xrightarrow{x_i^n} A\qq x_i)\notag\\
  &= \bigotimes_{i=1}^{r}(A\qq x_i\oplus A \qq x_i[1])\notag\\
  &= \bigoplus_{k=0}^{r} A_1^{\oplus \binom{r}{k} }[k].
\end{align}


Thus we learn that the terms in the tower $\{A_n\otimes_A A_1\}$ are independent of $n.$
Further from \ref{eq: formula for antimesa1} there is an obvious morphism of pro-systems $\{A_n\otimes_A A_1\}_{n>0}\to \{A_1\}_{n>0}$ corresponding to projection on the first factor and a section  $\{A_1\}_{n>0}\to \{A_n\otimes_A A_1\}_{n>0}.$ 

It remains to check that this induced morphism is actually a pro-isomorphism. For this it suffices to show that for any $n>1$ the morphism $A_n\otimes_A A_1\to A_{n-1}\otimes_A A_1$ is identity on the first factor of \ref{eq: formula for antimesa1} and $0$ on the other factors. This will prove the result as it is now clear that the morphism of pro-systems $\{A_n\otimes_A A_1\}_{n>0}\to \{A_1\}_{n>0}$ is an isomorphism. 

So now we prove the fact about the morphism being $\id$ on the first factor and $0$ on positive factors.

Note that the morphism $A_n\to A_{n-1}$ is obtained by base change along $\ZZ[x_1,\ldots, x_r]\xrightarrow{x_i\to f_i} A$ of the universal morphism 
\begin{equation}\label{eq: universal quotients}
  \ZZ[x_1,\ldots, x_r]/(x_1^n,\ldots, x_r^n)\to \ZZ[x_1,\ldots, x_r]/(x_1^{n-1},\ldots, x_r^{n-1})  
\end{equation}


The tensor product $A_n\otimes_A A_1\to A_{n-1}\otimes_A A_1$ is obtained by a further base change along $A\otimes_{\ZZ[x_1,\ldots, x_r]} \ZZ$. Thus to understand the morphism we may just base change \ref{eq: universal quotients} along $-\otimes_{\ZZ[x_1,\ldots, x_r]} \ZZ$.

Writing $$\ZZ[x_1,\ldots, x_r]/(x_1^n,\ldots, x_r^n)=\bigotimes_{i=1}^{i=r}\cofib(\ZZ[x_1,\ldots, x_r]\xrightarrow{x_i^n} \ZZ[x_1,\ldots, x_r])$$ and viewing the induced morphism of \ref{eq: universal quotients} by coming, after iterated tensoring, from the complex

$$% https://tikzcd.yichuanshen.de/#N4Igdg9gJgpgziAXAbVABwnAlgFyxMJZABgBpiBdUkANwEMAbAVxiRAB12AtL5ADwD6ARlKcGUCDjikABIIBOFEAF9S6TLnyEUIytXrNWiDt16CRYiVNkKlq9djwEiIoVVqMWbTj37DR7OKS0nICiipqIBiOWkRkbvqeRia+5gFB1qHhyu4wUADm8ESgAGbyEAC2SGQgOBBIIh6GbIJYAHqE9iBllQ3UdUgATInNxpxYUBGl5VWINQOIAMwjXsatU90zSMu19YjDTasgrW3AYAC0QsoqFMpAA
\begin{tikzcd}
{\ZZ[x_1,\ldots, x_r]} \arrow[r, "x_i^n"] \arrow[d, "x_i"] & {\ZZ[x_1,\ldots, x_r]} \arrow[d, "\id"] \\
{\ZZ[x_1,\ldots, x_r]} \arrow[r, "x_i^{n-1}"]              & {\ZZ[x_1,\ldots, x_r]}                 
\end{tikzcd},$$ we see that the derivation of \ref{eq: formula for antimesa1}, gives that the induces morphism $$\bigoplus_{k=0}^{r}\ZZ^{\oplus \binom{r}{k}}[k]\to \bigoplus_{k=0}^{r}\ZZ^{\oplus \binom{r}{k}}[k]$$ is $\id$ on the first factor and zero on all factors with a postive degree. After base change to $A_1$ we conclude the same for the induced morphism  

$$\bigoplus_{k=0}^{r} A_1^{\oplus \binom{r}{k} }[k]\to \bigoplus_{k=0}^{r} A_1^{\oplus \binom{r}{k} }[k].$$

This concludes the proof.


\end{proof}

\end{document}
